\chapter{Organismos, autoridades y responsabilidades}\label{ch:13}
\objetivos{Diferenciar normas internacionales, derecho de la Unión, supervisión nacional y datos del fabricante; localizar la autoridad competente para una pregunta concreta.}
\section{Un sistema con varios niveles}
La aviación internacional necesita compatibilidad entre Estados, pero ello no implica una autoridad única que apruebe todas las aeronaves y todas las tareas. La OACI desarrolla normas y métodos recomendados internacionales dentro del sistema del Convenio de Chicago. Los Estados asumen responsabilidades de implementación y supervisión según sus obligaciones y marco jurídico. Una norma internacional y una instrucción de mantenimiento de un motor no tienen la misma función.

En la Unión Europea, el Reglamento 2018/1139 establece reglas comunes y la Agencia de la Unión Europea para la Seguridad Aérea (EASA). Los reglamentos de desarrollo distribuyen requisitos técnicos y administrativos. Las autoridades nacionales, como AESA en España, ejercen las competencias que les corresponden dentro del sistema. Debe comprobarse quién es competente en cada procedimiento y territorio.
\section{EASA y AESA}
EASA participa, entre otras funciones, en certificación de productos y elaboración del marco técnico; AESA realiza actuaciones de supervisión y aprobación dentro de sus competencias. No conviene memorizar «EASA diseña y AESA mantiene»: ni una agencia diseña el producto del fabricante ni esa frase representa la distribución jurídica completa.

La autoridad competente de una organización depende de las reglas aplicables, localización y posibles reasignaciones. Por ejemplo, 145.1 identifica la autoridad para Parte 145; no se infiere solo de que la aeronave vuele a España. La aprobación de una CAMO es diferente de la certificación de un motor y de la licencia de una persona.
\section{FAA y el sistema estadounidense}
La FAA es la autoridad de aviación civil de Estados Unidos dentro de su sistema legal. Las disposiciones del título 14 del CFR y las directivas FAA deben interpretarse en ese ámbito. Un documento estadounidense puede ser una fuente técnica excelente sin convertirse automáticamente en requisito europeo.

La aplicabilidad en una aeronave concreta necesita considerar Estado de matrícula, Estado de diseño, adopción o reconocimiento de requisitos y configuración. Una AD FAA relativa a un motor puede ser relevante para ese producto instalado en múltiples modelos, pero no basta con encontrar el nombre comercial en una hoja para cerrar su aplicabilidad en cualquier jurisdicción.
\section{Fabricante, titular de diseño y operador}
El titular de un certificado de tipo mantiene responsabilidades relacionadas con el diseño y la aeronavegabilidad continuada; los titulares de cambios o reparaciones también aportan datos dentro de su ámbito. El operador y el responsable de la aeronavegabilidad gestionan obligaciones asociadas a la explotación y al estado de la aeronave. La organización de mantenimiento realiza tareas dentro de su aprobación y procedimientos. El personal certificador ejerce privilegios y autorizaciones concretos.

Estas funciones pueden coexistir en una misma entidad, pero no pierden su diferencia. Una organización que diseña una reparación no queda por ello autorizada a ejecutar cualquier mantenimiento; una CAMO no debe confundirse con el taller que realiza físicamente todas las tareas.
\begin{table}[htbp]\centering\small\caption{Preguntas de competencia y fuente de partida.}\label{tab:autoridades}
\begin{tabular}{p{.43\linewidth}p{.45\linewidth}}\toprule
Pregunta & Punto de partida\\\midrule
¿Está el producto dentro del sistema europeo? & Reglamento 2018/1139, artículo 2 y anexos relevantes\\[4pt]
¿Quién supervisa una Parte 145? & Reglamento 1321/2014, 145.1\\[4pt]
¿Quién aporta instrucciones de aeronavegabilidad continuada? & Reglamento 748/2012, 21.A.7 y datos del titular aplicable\\[4pt]
¿Quién gestiona aeronavegabilidad y bajo qué régimen? & Reglamento 1321/2014, artículo 3, M.A.201 o ML.A.201 y parte organizativa aplicable\\\bottomrule
\end{tabular}\end{table}
\section{Cómo distinguir obligación y orientación}
Los reglamentos contienen obligaciones jurídicas dentro de su ámbito. Los AMC describen medios aceptables de cumplimiento; los GM ofrecen orientación. Las CS forman parte del sistema de especificaciones usado para establecer bases de certificación; la base aplicable a un producto no es necesariamente la última enmienda publicada.

Un boletín de servicio puede recomendar una acción y una AD puede imponer acciones que incorporan un documento por referencia. Las palabras «mandatory» o «alert» en un documento de fabricante no resuelven por sí solas todas las vías por las que una tarea puede ser obligatoria. Hay que rastrear AD, limitaciones, programa y reglas aplicables.
\caso{M13. Una afirmación sin ámbito}{«Este curso permite certificar motores» es una conclusión injustificada si solo se acredita una formación universitaria. Se pregunta qué acto se pretende certificar, qué licencia o aprobación se exige, qué habilitación existe y qué autorización emite la organización. La titulación puede aportar conocimientos y ser relevante para determinadas rutas, pero no sustituye esos requisitos.}
\section{Método de consulta}
Toda pregunta normativa se transforma en una ficha: objeto, Estado y operación, autoridad, norma y versión, disposición, fecha de aplicación y conclusión limitada. Si falta una de esas piezas, se conserva la incertidumbre y se busca el documento correspondiente. Una página de orientación no debe usarse para contradecir sin análisis un texto vinculante posterior.

\section{Del marco internacional al expediente concreto}
OACI proporciona el marco internacional y normas y métodos recomendados dentro de su sistema; su existencia no convierte cada documento en una orden ejecutable por un taller. En la Unión Europea, el Reglamento básico y sus actos de desarrollo distribuyen funciones y requisitos. Para un problema se identifica jurisdicción, producto, matrícula, operación y actividad antes de escoger la disposición.

Una autoridad puede aprobar o supervisar sin ser autora de todas las instrucciones técnicas. El titular de diseño emite información dentro de su función; una organización desarrolla procedimientos dentro de su aprobación. Distinguir emisor, aprobador, usuario y supervisor evita atribuir a un manual funciones jurídicas que no tiene.
\section{Obligación, medio de cumplimiento y orientación}
Un requisito vinculante expresa qué debe satisfacerse dentro de su ámbito. Un AMC ofrece una vía aceptable de demostrar cumplimiento del requisito asociado; un GM ayuda a interpretarlo. Los medios alternativos tienen procedimientos específicos. No toda frase de un compendio EAR tiene la misma fuerza por aparecer en el mismo PDF.

Una especificación de certificación se relaciona con una base y un procedimiento. Una nueva enmienda no sustituye automáticamente toda base anterior. Un SB procede del fabricante y puede resultar obligatorio por incorporación en otro requisito o dato vinculante; su título por sí solo no resuelve la cuestión.
\section{Distribuir responsabilidades sin perder trazabilidad}
Un expediente académico indica quién conserva programa, planifica, ejecuta, certifica y revisa. No asigna todo a «EASA» ni todo al mecánico. La contratación de una tarea no elimina por sí sola todas las responsabilidades del contratante. La estructura exacta depende de régimen, operación y organizaciones.

La hoja de control es una vista del sistema documental; no crea una aprobación. Si un registro es inconsistente se conserva evidencia y se activa una ruta de aclaración y corrección. Un cálculo interno correcto no sustituye la demostración de ejecución exigible.

\section{Convenio de Chicago, IATA y el antecedente JAA}
El Convenio de Chicago establece el marco internacional del que forman parte los estándares y métodos recomendados de OACI. Esos instrumentos buscan armonización entre Estados; su relación con la norma aplicable se analiza a través de los compromisos y sistemas pertinentes. No son una orden de mantenimiento de un componente ni deben confundirse con el manual del fabricante. La explicación institucional de OACI sobre elaboración de SARPs ayuda a identificar esa función \cite{ICAO-SARPS}.

IATA es una asociación de compañías aéreas. Su actividad sectorial, estándares y programas pueden ser relevantes para operadores, pero no la convierten en autoridad que expide licencias Parte 66 o aprueba una reparación en nombre de AESA. Se identifica la naturaleza del instrumento y, cuando se usa, su relación contractual, técnica o regulatoria concreta \cite{IATA}.

Las JAA son un antecedente del sistema europeo actual de armonización y sus JAR. EASA explica la transición desde una cooperación de autoridades nacionales hacia el marco común. Una referencia histórica JAR-145 en un libro antiguo no se usa como si sustituyera la Parte 145 vigente. Se conserva su utilidad histórica y se verifica el texto actual para cada afirmación \cite{EASA-JAA}.

\repaso{¿Una guía publicada en una web oficial es necesariamente legislación?}{La procedencia identifica la fuente, pero se distingue naturaleza y alcance. La guía se usa junto al requisito, su versión y sus condiciones.}

\resumen{La competencia se determina por ámbito y procedimiento. OACI, EASA, AESA, FAA y fabricante no cumplen funciones intercambiables. La etiqueta de un documento no sustituye el análisis de su fuerza y aplicabilidad. Una formación académica no equivale a una autorización de certificación.}
\fuentes{Reglamento 2018/1139, arts. 2, 62 y 75--77 \cite{EU1139}; 1321/2014, 145.1 \cite{EU1321}; AESA, procedimiento CAMO \cite{AESA-CAMO}; materiales docentes de marco normativo.}
